无限均匀、无碰撞的冷电子等离子体处于沿 \(+z\) 的恒定外磁场 \(\mathbf B_0=B_0\hat{\mathbf z}\) 中，离子视为静止中性背景，电子数密度为 \(n_e>0\)、质量为 \(m_e\)、电荷为 \(-e\)（\(e>0\)）。忽略热压、辐射反作用和非线性响应。定义电子等离子体频率 \(\omega_p^2=n_e e^2/(\epsilon_0m_e)\)、回旋频率 \(\omega_c=eB_0/m_e>0\)。一束沿 \(+z\) 传播的横向单色电磁波采用 \(e^{i(kz-\omega t)}\) 复振幅约定，\(c=1/\sqrt{\epsilon_0\mu_0}\)。定义圆偏振分量 \(E_\pm=E_x\pm iE_y\)，仅考虑实正折射率对应的传播支。
(1) 从电子线性化洛伦兹力方程和 Maxwell 方程出发，求 \(v_\pm,J_\pm\) 与两个圆偏振本征折射率 \(n_\pm^2(\omega)\)，说明符号与 \(e^{-i\omega t}\) 约定的对应关系。
(2) 分析两支传播模的截止与电子回旋共振：分别求正频率截止 \(\omega_{\mathrm{cut},+}\)、\(\omega_{\mathrm{cut},-}\)，指出 \(\omega=\omega_c\) 附近 \(n_-^2\) 的符号变化与无碰撞模型的适用限制。
(3) 令入射线偏振光在 \(z=0\) 沿 \(x\) 方向，频率高于两支截止且高于 \(\omega_c\)。经过长度 \(L>0\) 的等离子体后，求偏振面旋转角 \(\Theta\)（按 \(x\) 朝 \(y\) 为正），推导 \(\omega\gg\omega_p,\omega_c\) 的 Faraday 旋转近似及其频率标度。
(4) 推导两圆偏振模的群时延 \(t_{g,\pm}=L\,dk_\pm/d\omega\) 的闭式表达式，证明差分群时延与旋转角频率导数的精确关系；求高频下的首项差分群时延，并检查零磁场与零密度极限。
